\documentclass[11pt]{article}
\usepackage[margin=1in]{geometry}
\usepackage{amsmath,amssymb,mathtools}
\usepackage{booktabs}
\usepackage{enumitem}
\usepackage{siunitx}
\sisetup{group-separator={,},group-minimum-digits=3}

\title{\textbf{Solutions -- Quiz 2} \vspace{ 0.4cm }\\ECON 300: Labour Economics}
\author{Instructor: Muhebullah Karimzada}
\date{February 05,  2026}

\begin{document}
\maketitle




\subsection*{Question 1}

\textbf{Correct answer: (c)}

The slope \( (1-0.6)W \) represents the \emph{effective net wage} faced by the worker while EI benefits are being reduced.  
For each additional dollar earned, EI benefits are reduced by 60 cents, so disposable income rises by only 40 cents.  
This implicit reduction in earnings is equivalent to a marginal tax created by the EI program.



\subsection*{Question 2}

\textbf{Correct answer: (b)}

Between points \textbf{B} and \textbf{D}, EI benefits have been fully exhausted.  
As a result, additional work no longer reduces benefits, and the worker keeps the full market wage \(W\).  
This is why the slope of the budget constraint in this region equals \(W\).



\subsection*{Question 3}

\textbf{Correct answer: (c)}

Point \textbf{A} corresponds to zero or minimal weeks worked with full EI benefits received.  
At this point, the worker enjoys the maximum benefit level provided by the EI system while supplying little or no labour.  
This point represents the income support available at the lowest level of work.



\subsection*{Question 4}

\textbf{Correct answer: (c)}

The budget constraint is flatter between \textbf{A} and \textbf{B} because EI benefits are clawed back as earnings increase.  
Each additional unit of work leads to a reduction in benefits, lowering the net return to working.  
This flattening reflects the implicit tax imposed by the benefit reduction rate.



\subsection*{Question 5}

\textbf{Correct answer: (c)}

The difference between the gross wage \(W\) and the effective wage \( (1-0.6)W \) reflects the implicit tax created by the EI benefit reduction rate.  
This is not a statutory income tax or payroll tax, but rather a reduction in benefits that lowers disposable income from additional work.



\subsection*{Question 6}

\textbf{Correct answer: (c)}

The highest effective marginal tax rate occurs between points \textbf{A} and \textbf{B}.  
In this region, each dollar of earnings reduces EI benefits by 60 cents, implying a high effective marginal tax rate.  
Once benefits are exhausted (after point \textbf{B}), the marginal tax rate from EI falls to zero.



\subsection*{Question 7}

\textbf{Correct answer: (b)}

Compared to a world without EI, the EI scheme increases income when not working or working little.  
This can reduce labour force participation for some individuals by lowering the incentive to enter employment.  
This effect operates primarily on the extensive margin of labour supply.



\subsection*{Question 8}

\textbf{Correct answer: (c)}

Point \textbf{B} is a kink in the budget constraint where the slope changes discretely.  
At such kinks, workers may optimally locate because small changes in labour supply lead to large changes in marginal incentives.  
With standard convex preferences, individuals may bunch at these kink points.



\subsection*{Question 9}

\textbf{Correct answer: (c)}

If the EI replacement rate increases from 0.6 to 0.8, the net-of-benefit-reduction share falls from 0.4 to 0.2.  
This makes the slope of the EI segment smaller, meaning the budget constraint becomes flatter while on EI.  
A flatter constraint implies weaker incentives to increase work effort.



\subsection*{Question 10}

\textbf{Correct answer: (b)}

The key efficiency concern illustrated by the diagram is that EI creates high effective marginal tax rates.  
These high marginal rates discourage work by reducing the net return to additional labour.  
This disincentive effect is central to policy debates surrounding unemployment insurance design.



\section*{Summary of Correct Answers}

\begin{center}
\begin{tabular}{cccccccccc}
1 & 2 & 3 & 4 & 5 & 6 & 7 & 8 & 9 & 10 \\
\hline
C & B & C & C & C & C & B & C & C & B
\end{tabular}
\end{center}

\vfill
\textbf{End of Solutions}

\end{document}
